\documentclass[10pt,a4paper]{amsart}
\usepackage[margin=19mm]{geometry}
\usepackage{amsmath,amssymb,mathtools}
\usepackage[scr=rsfs,cal=euler]{mathalfa}
\usepackage{xfrac}
\usepackage[unicode,hidelinks,bookmarks=false]{hyperref}
\newcommand{\ie}{i.e.\ }
\newcommand{\cf}{cf.\ }
\newcommand{\resp}{respectively\ }
\newcommand{\Subsection}{\subsection}
\providecommand{\nobreakdash}{\nobreak\hbox{-}}
\setlength{\emergencystretch}{2em}
\multlinegap0pt
\usepackage{bm}
\usepackage{bbm}
\usepackage{dsfont}
\usepackage{xspace}
\usepackage{cancel}
\usepackage{mathtools}
\usepackage{amsthm}
\makeatletter
\@ifundefined{theorem}{\newtheorem{theorem}{Theorem}}{}
\@ifundefined{lem}{\newtheorem{lem}[theorem]{\sc Lemma}}{}
\makeatother
\title[Generalized torsion elements in groups]{Generalized torsion elements in groups}
\author{Raimundo Bastos, Csaba Schneider, Danilo Silveira}
\date{Source: arXiv:2302.09589v1 (CC BY 4.0)}
\begin{document}
\maketitle

\noindent\emph{Source and licence.} Generalized torsion elements in groups. Version arXiv:2302.09589v1 (CC BY 4.0), distributed under the Creative Commons Attribution 4.0 International licence, as recorded in the arXiv licence metadata for this exact version and shown on the abstract page https://arxiv.org/abs/2302.09589v1. Full rights notice: licenses/arxiv-2302.09589-CC-BY-4.0.txt.\par\smallskip

\noindent\textit{Editorial scope.} Counted unit: the Lem whose source label is \texttt{lem:commutators-powers} in the article (source0.tex lines 347--360, source label \texttt{lem:commutators-powers}), delivered together with the article's own complete original proof of it (source0.tex lines 361--399, one \texttt{proof} environment of 39 lines). No mathematics is written, repaired or re-derived: statement and proof are the article's own text line for line. Required context retained uncounted: prop\_comm (lines 333--342). Every definition, display and label of those lines is retained unchanged. Omitted from this package and not redistributed: 1--332, 343--346, 400--535. Nothing inside a retained excerpt is omitted or altered: every retained body is delivered exactly as the article's lines are written, so no omission receipt of any kind accompanies this package. Citations: the retained excerpts carry no citation command at all, so no bibliography entry of the article is transcribed and no reference is redistributed. Reference adaptation: none was needed and none was made; every reference inside the delivered unit resolves inside it, so no unresolved reference or citation is produced. Printed slips kept literal: no printed word, symbol, hypothesis or number of the counted statement or of its proof was altered. Layout and class: the article's own class is \texttt{amsart} with packages including amssymb, amsthm, amsmath, latexsym, wasysym, soul, xy, xcolor; the wrapper is the lane's pinned \texttt{amsart} editorial wrapper at 10pt on A4, which restates the theorem-like environments the retained text needs and adds the article's own macro definitions that the retained text needs (none), with extra wrapper packages (bm, bbm, dsfont, xspace, cancel, mathtools). The delivered numbering is the unit's own. Assets: no retained span contains a figure environment, a graphics-inclusion command, a PDF-page inclusion, a table or a TikZ picture; no image file is copied into this package, the package declares no asset dependency and no image file is redistributed with it. Licence: version arXiv:2302.09589v1 of this article is distributed under the Creative Commons Attribution 4.0 International licence, as recorded in the arXiv licence metadata for this exact version and shown on the abstract page https://arxiv.org/abs/2302.09589v1. A full rights notice is delivered at licenses/arxiv-2302.09589-CC-BY-4.0.txt. Characters: every retained excerpt is pure ASCII, so the delivered engine, XeLaTeX with its default Latin Modern text font, reports no missing character.
\begin{lem}\label{prop_comm}
For elements $x, y, z$ of a group $G$ and a positive integer $k$, the following identities are valid:
\begin{enumerate}
\item $xy=yx[x,y]$
\item $x^y=x[x,y]$
    \item $[xy,z]=[x,z][x,z,y][y,z]$
    \item $[x^k,y]=[x,y]^{x^{k-1}}[x,y]^{x^{k-2}}\cdots [x,y]^x[x,y]$
   \item $x^ky^k=(xy)^kc_2^{\binom{k}{2}}\ldots c_i^{\binom{k}{i}}\ldots c_{k-1}^kc_k$, where $c_i\in \gamma_i(\left<x,y\right>)$ for each non-negative integer $i$.
\end{enumerate}
\end{lem}

\begin{lem} \label{lem:commutators-powers}
Let $k\geq 2$, and let $x, g_1,\ldots,g_k$ be elements in a group $G$.
\begin{enumerate}
    \item We have that  
    \[
    x^{g_1}x^{g_2}\cdots x^{g_k}=x^k\sigma_2
    \]
where $\sigma_2$ is a product of commutators of weight at least $2$ and the element $x$ is the first entry of all the factors of $\sigma_2$.
    \item If $o_{\bullet}(x)=k$, then  $x^{k}=c_1\cdots c_r$, where each $c_i=[c_{i,1},c_{i,2}]$ is a commutator of weight at least $2$ such that the first entry of $c_{i,2}$ is $x$.
    \item If $o_{\bullet}(x)=k$ and $c_m$ is a commutator of weight $m$ with $x$ in some entry, then $c_m^k$ is a product of commutators of weight at least $m+1$  and $x$ appears in some entry of all factors of $c_m^k$.
    \item If $o_{\bullet}(x)=k$ and $\sigma_m$ is a product of commutators of weight at least $m$ with $x$ in some entry of all its factors, then $\sigma_m^k$ is a product of commutators of weight at least $m+1$  and $x$ appears in some entry of all its factors.
    \item If $o_{\bullet}(x)=k$, then $(x^{g_1}x^{g_2}\cdots x^{g_k})^{k^m}=x^{k^{m+1}}\sigma_{m+1}$, where $\sigma_{m+1}$ is a product of commutators of weight at least $m+1$ and the element $x$ appears in some entry of all factors of $\sigma_{m+1}$.
  \end{enumerate}
\end{lem}

\begin{proof}
(1) We proceed by induction on $k.$ If $k=2$, then 
$$x^{g_1}x^{g_2}=x[x,g_1]x[x,g_2]=x^2[x,g_1][x,g_1,x][x,g_2].$$ Assuming the result holds for $k\geq 2,$ we get 
$$(x^{g_1}x^{g_2}\cdots x^{g_k})x^{g_{k+1}}=(x^k\sigma_2)x[x,g_{k+1}]=x^{k+1}\sigma_2[\sigma_2,x][x,g_{k+1}],$$ where $\sigma_2$ is a product of commutators of weight at least $2$ and the element $x$ is the first entry of the factors of $\sigma_2$. 
Now, the result follows by applying the Lemma \ref{prop_comm}(3) several times to the commutator $[\sigma_2,x]$.

(2)  If $o_{\bullet}(x)=k$, then there exist elements $g_1,\ldots, g_k\in G$ such that   $1=x^{g_1}x^{g_2}\cdots x^{g_k}$. By previous item, we can write $$1=x^{g_1}x^{g_2}\cdots x^{g_k}=x^k\sigma_2,$$ 
where $\sigma_2$ is a product of commutators of weight at least $2$ and the element $x$ is the first entry of all the factors of $\sigma_2$. Thus, $x^k=\sigma_2^{-1}$.

(3)  We proceed by induction on $m$. The basic step $m=1$ follows by item (2). Assume the result holds for all positive integers up to $m$. Let $c_{m+1}$ be a commutator of weight $m+1$ with $x$ in some entry.  Write $c_{m+1}=[c_i,c_j]$ where $c_i, c_j$ are commutators of weight $i$ and $j$, respectively, and $i+j=m+1.$ Without loss of generality, we can assume that $x$ occurs in the left sub-commutator $c_i$. By Lemma \ref{prop_comm}(4) we get
\begin{align*}
   [c_i^k,c_j]& =    [c_i,c_j]^{c_i^{k-1}}[c_i,c_j]^{c_i^{k-2}}\cdots [c_i,c_j]^{c_i}[c_i,c_j]\\
    & =    ([c_i,c_j][c_i,c_j,c_i^{k-1}])
    %([c_i,c_j][c_i,c_j,c_i^{k-2}])
    \cdots ([c_i,c_j][c_i,c_j,c_i])[c_i,c_j]\\
    & =  [c_i,c_j]^k\sigma_{m+2}
\end{align*}
where $\sigma_{m+2}$ is a product of commutators of weight at least $m+2$ with $x$ in some entry of all factors. By the induction hypothesis, $c_i^k$ is a product of commutators of weight at least $i+1$  and $x$ appears in some entry of all its factors, say $c_i^k=c_{i,1}\cdots c_{i,r}$. Since 
\[
c_{m+1}^k=[c_i,c_j]^k=[c_i^k,c_j]\sigma_{m+2}^{-1},
\]
the result follows by applying several times Lemma \ref{prop_comm}(3) to $[c_i^k,c_j]=[c_{i,1}\cdots c_{i,r},c_j]$.


(4) We proceed by induction on the number $r$ of factors of $\sigma_m$. The basic step $r=1$ follows by the previous item. Assume that the result holds for all positive integers up to $r\geq 1$ and set $\sigma_m=\tau_1\cdots\tau_r\tau_{r+1}$ where each $\tau_i$ is a commutator of weight at least $m$ with $x$ in some entry. By Lemma \ref{prop_comm}(5) we can deduce that 
$$\sigma_m^k=(\tau_1\cdots\tau_r\tau_{r+1})^k=(\tau_1\cdots\tau_r)^k\tau_{r+1}^k\sigma_{m+1}$$
where $\sigma_{m+1}$ is a product of commutators of weight at least $m+1$ with $x$ appearing in some entry in each factor. Thus, the result follows applying the induction hypothesis to $(\tau_1\cdots\tau_r)^k$ and to $\tau_{r+1}^k$. 

(5)  We show item (5) by induction on $m$. Firstly, we will show the basic step $m=1$. By item (1), we can write 
$x^{g_1}x^{g_2}\cdots x^{g_k}=x^k\sigma_2$. 
By parts (1) and (5) of Lemma \ref{prop_comm}, we have  
$$(x^{g_1}x^{g_2}\cdots x^{g_k})^k=(x^k\sigma_2)^k=x^{k^2}\tilde{\sigma_2}$$
where $\tilde{\sigma_2}$ is a product of commutators of weight at least 2 and the elements $x$ appears in some entry of all factors.

 Now, assume that the result holds for  all positive integers up to $m\geq 1$. By the induction hypothesis, we  get  
$$(x^{g_1}x^{g_2}\cdots x^{g_k})^{k^{m+1}}=((x^{g_1}x^{g_2}\cdots x^{g_k})^{k^{m}})^k=(x^{k^{m+1}}\sigma_{m+1})^k$$ 
%
where $\sigma_{m+1}$ is a product of commutators of weight at least $m+1$ and the element $x$ appears in some entry of all factors. By Lemma \ref{prop_comm}(5), we can deduce that $(x^{k^{m+1}}\sigma_{m+1})^k=x^{k^{m+2}}\sigma_{m+1}^k\sigma_{m+2}$, where $\sigma_{m+2}$ is a product of commutators of weight at least $m+2$ and the element $x$ appears in some entry of all the factors. Thus, the result follows applying item (4) to $\sigma_{m+1}^k$.
\end{proof}

\end{document}
